problem:  
Let \( M^m \subset \mathbb{S}^{m+k}(1) \) be a compact properly biharmonic submanifold with constant mean curvature vector \(H\).  
It is known that for biharmonic hypersurfaces (\(k=1\)) of \(\mathbb{S}^{m+1}\) with constant mean curvature, all examples must satisfy \(\|H\|^2 = \tfrac12\); the canonical examples include the small hypersphere \(\mathbb{S}^m(1/\sqrt{2})\) and certain generalized Clifford tori. These results are dimensionally rigid in codimension one.  

**Conjecture (Higher-Codimensional Gap Phenomenon):**  
There exists a constant \(\delta_m > 0\), depending only on the dimension \(m\) of \(M\) but not on the codimension \(k\), such that for every compact, properly biharmonic, constant-mean-curvature submanifold \(M^m \subset \mathbb{S}^{m+k}\), one has either  
\[
\|H\|^2 = \tfrac{1}{2} \quad \text{or} \quad \|H\|^2 \ge \tfrac{1}{2} + \delta_m.
\]
In particular, there should be no sequence of properly biharmonic submanifolds in spheres whose mean curvature values accumulate to \(\tfrac{1}{2}\) from above.  

A further question: Is the existence of such a universal \(\delta_m\) implied by any geometric or analytical constraint — e.g., parallel mean curvature vector field, flat normal connection, or positivity of the normal Ricci curvature — and can it be derived from a Simons-type inequality associated with the biharmonic equation?  

Why is it a "good" problem:  
This formulation correctly rests on the known classification of constant-mean-curvature biharmonic hypersurfaces and transforms that knowledge into a sharply posed conjecture probing whether a *gap phenomenon* persists in higher codimensions.  
It is a “good” problem for several reasons:  

- **Depth and Clarity:** The statement is simple — does a universal gap \( \delta_m \) exist for \(\|H\|^2\)? — yet addressing it demands new geometric and analytic insight into the biharmonic equation in higher codimension.  
- **Bridge Across Fields:** It bridges geometric analysis, spectral theory, and submanifold geometry, since any proof will likely require extending Simons-type identities, integral estimates, or curvature pinching techniques from the harmonic to the biharmonic setting.  
- **Predictive Power:** If true, it would reveal a quantization-like rigidity of biharmonic energy in spheres, suggesting deep structural restrictions on curvature beyond the currently known codimension-one cases.  
- **Fertile Ground:** Even partial results (e.g., establishing the gap under additional geometric conditions such as parallel mean curvature) would significantly advance the understanding of biharmonic geometry and might lead toward a full classification in higher codimension.  

Thus, the problem stands at the intersection of known rigidity phenomena and the unexplored analytic subtleties of the biharmonic map equation, representing a clean, profound, and open research challenge.